\documentclass[11pt,a4paper]{article}

% ---------- Packages ----------
\usepackage[margin=1in]{geometry}
\usepackage{amsmath, amssymb, amsthm, mathtools}
\usepackage{microtype}
\usepackage[colorlinks=true,
            linkcolor=blue!50!black,
            urlcolor=blue!50!black,
            citecolor=blue!50!black]{hyperref}

% ---------- Theorem environments ----------
\theoremstyle{plain}
\newtheorem{theorem}{Theorem}[section]
\newtheorem{proposition}[theorem]{Proposition}
\newtheorem{lemma}[theorem]{Lemma}
\newtheorem{corollary}[theorem]{Corollary}

\theoremstyle{definition}
\newtheorem{definition}[theorem]{Definition}

\theoremstyle{remark}
\newtheorem{remark}[theorem]{Remark}
\newtheorem{example}[theorem]{Example}

% ---------- Math shortcuts ----------
\newcommand{\R}{\mathbb{R}}
\newcommand{\N}{\mathbb{N}}
\newcommand{\E}{\mathbb{E}}
\newcommand{\Var}{\operatorname{Var}}
\newcommand{\F}{\mathcal{F}}
\newcommand{\Prob}{\mathbb{P}}
\newcommand{\Q}{\mathbb{Q}}
\newcommand{\dd}{\mathop{}\!\mathrm{d}}
\newcommand{\eqdef}{\overset{\mathrm{def}}{=}}
\newcommand{\Lop}{\mathcal{L}}

% ---------- Document metadata ----------
\title{SDE Discretisation:\\
       It\^o-Taylor Expansions, Euler--Maruyama, and Milstein Schemes}
\author{Quant Finance Portfolio --- Phase 2, Block 1, Algorithm 2 (theory)}
\date{\today}

\begin{document}
\maketitle

\begin{abstract}
The exact GBM sampler used in Block~1.1 exploited a closed-form
solution of the underlying SDE. Beyond Black--Scholes, closed forms
are essentially unavailable: pricing under local volatility, stochastic
volatility, jump-diffusion, or path-dependent payoffs all require
\emph{simulating} approximate trajectories of the underlying
stochastic process by discretising the SDE in time. This document
develops the foundations of that discretisation: the two notions of
convergence relevant to SDEs (\emph{strong} and \emph{weak}), the
It\^o-Taylor expansion that generates discretisation schemes
systematically, and the two canonical schemes of order one in the
relevant senses (Euler--Maruyama and Milstein). The emphasis is
conceptual: the precise statements of convergence theorems are given
without proof (Kloeden--Platen is the canonical reference), but the
\emph{reason} behind each convergence order is developed in detail
through the It\^o-Taylor expansion, and the practical consequences for
finance applications are spelled out.
\end{abstract}

\tableofcontents

\section{Introduction}
\label{sec:intro}

Block~1.1 priced the European call by drawing $S_T$ from its exact
distribution under geometric Brownian motion. The construction was
trivial because the GBM SDE
\begin{equation}
\dd S_t \;=\; r S_t \dd t + \sigma S_t \dd W_t
\end{equation}
admits a closed-form solution, $S_T = S_0 \exp((r - \sigma^2/2)T +
\sigma W_T)$, so a single normal draw suffices to generate one path's
terminal value. The exactness was a luxury, not a method.

Almost every model of practical interest beyond Black--Scholes lacks
such a closed form. Local volatility, where $\sigma$ depends on $t$
and $S_t$, has no closed form for $S_T$ even in simple parameterisations.
Stochastic volatility models like Heston have a known characteristic
function but no closed form for the joint $(S_T, V_T)$.
Jump-diffusion models add a Poisson component that destroys the
log-normal structure entirely. In all these cases, the only practical
route to $\E[f(S_T)]$ is to simulate \emph{approximate} trajectories
by discretising the SDE in time and then averaging.

The discretisation introduces a new kind of error, qualitatively
different from the statistical error of Block~1.1. The Monte Carlo
estimator was unbiased: any deviation between $\hat C_n$ and
$C_{\mathrm{BS}}$ was purely sampling error, controlled by $1/\sqrt{n}$.
Once we replace the exact $S_T$ by a discretised $\hat S_T$, the
estimator becomes \emph{biased}: even with $n = \infty$ paths, we
estimate $\E[f(\hat S_T)]$, not $\E[f(S_T)]$. The size of the bias
depends on the time step $h$ and on the scheme, and the central
analytical question of this block is: at what rate does this bias
vanish as $h \to 0$?

The answer is not a single number. SDE solutions admit two distinct
and largely independent notions of approximation, called \emph{strong}
and \emph{weak} convergence. The strong notion measures pathwise
proximity, the weak notion measures distributional proximity, and a
discretisation scheme can be of different order in the two senses.
The most consequential example is the Euler--Maruyama scheme, which
has strong order $\tfrac{1}{2}$ but weak order $1$.

The remainder of this document is structured as follows.
Section~\ref{sec:setting} fixes the SDE setting and introduces the
abstract notion of a discretisation scheme.
Section~\ref{sec:convergence} develops the two notions of convergence
in parallel, with examples and an explicit relation between them.
Section~\ref{sec:ito-taylor} derives the It\^o-Taylor expansion at the
level of detail required to read off discretisation schemes from it.
Section~\ref{sec:euler} derives the Euler--Maruyama scheme as a
first-order truncation and states its convergence orders.
Section~\ref{sec:milstein} does the same for the Milstein scheme as a
second-order truncation. Section~\ref{sec:practical} discusses which
scheme is appropriate in which context. Section~\ref{sec:roadmap}
sets up the implementation work of Blocks~1.2.1 and~1.2.2.

\section{Setting and discretisation schemes}
\label{sec:setting}

\subsection{The scalar It\^o SDE}

We work with a scalar It\^o SDE
\begin{equation}
\label{eq:sde}
\dd X_t \;=\; a(X_t) \dd t \;+\; b(X_t) \dd W_t,
\qquad X_0 = x_0,
\qquad t \in [0, T],
\end{equation}
on a filtered probability space $(\Omega, \F, \{\F_t\}_{t \in [0,T]},
\Prob)$ supporting a standard Brownian motion $W$. The coefficients
$a, b: \R \to \R$ are taken sufficiently regular (Lipschitz, with
linear growth) for a unique strong solution to exist; we will state
the additional regularity required by each convergence theorem when
needed.

Geometric Brownian motion is the running example: $a(x) = r x$,
$b(x) = \sigma x$. We will refer back to it both for sanity checks and
because the closed-form solution of GBM provides the reference path
against which the discretised paths of Blocks~1.2.1 and~1.2.2 will be
validated empirically.

\subsection{Discretisation schemes}

Fix a partition
\begin{equation}
0 = t_0 < t_1 < \dots < t_N = T,
\qquad t_n = n h,
\qquad h = T / N,
\end{equation}
and the corresponding Brownian increments
\begin{equation}
\Delta W_n \;\eqdef\; W_{t_{n+1}} - W_{t_n}
\;\sim\; \mathcal{N}(0, h),
\qquad n = 0, 1, \dots, N-1.
\end{equation}
The increments $\Delta W_n$ are independent across $n$ by the
independent-increments property of Brownian motion.

A \emph{discretisation scheme} for~\eqref{eq:sde} is a deterministic
recursion of the form
\begin{equation}
\label{eq:scheme-abstract}
\hat X_{n+1} \;=\; \Psi\!\left(
\hat X_n,\ h,\ \Delta W_n,\ \dots
\right),
\qquad \hat X_0 = x_0,
\end{equation}
where $\Psi$ is some function of the current state, the step size, the
Brownian increment, and possibly higher-order random variables (we
will see exactly which in the It\^o-Taylor section). The output is a
sequence $\hat X_0, \hat X_1, \dots, \hat X_N$ that we interpret as
approximations to $X_{t_0}, X_{t_1}, \dots, X_{t_N}$.

Two clarifications. First, the partition is uniform throughout this
document. Adaptive step sizes exist and have their place but are not
needed here. Second, the scheme is deterministic given
$\Delta W_0, \dots, \Delta W_{N-1}$: same Brownian path, same scheme,
same output, no internal randomness in $\Psi$ beyond the increments.

\section{Strong and weak convergence}
\label{sec:convergence}

The fundamental distinction of SDE numerical analysis, absent from
the deterministic ODE setting, is that there is more than one natural
sense in which $\hat X_N$ approximates $X_T$. We develop the two
relevant notions and explore their relationship.

\subsection{Pathwise vs distributional approximation}

Suppose we run the scheme with the \emph{same} Brownian increments
$\Delta W_n$ that were used to build the exact solution $X_T$ on a
common probability space. Two natural questions about the quality of
the approximation suggest themselves:

\begin{enumerate}
\item How close is $\hat X_N$ to $X_T$, \emph{path by path}? That is,
on a typical realisation $\omega \in \Omega$, how large is
$|X_T(\omega) - \hat X_N(\omega)|$?
\item How close is the distribution of $\hat X_N$ to that of $X_T$?
That is, for nice test functions $f$, how large is
$|\E[f(\hat X_N)] - \E[f(X_T)]|$?
\end{enumerate}

These questions are formally distinct. The first requires that
$X_T$ and $\hat X_N$ live on the same probability space, which we
arrange by sharing the Brownian path. The second is about marginal
distributions and does not require any coupling: even if $\hat X_N$
were generated from an entirely separate Brownian motion, the
expectation $\E[f(\hat X_N)]$ is well-defined and the comparison with
$\E[f(X_T)]$ remains sensible.

The first question is what \emph{strong} convergence formalises; the
second, what \emph{weak} convergence formalises.

\subsection{Strong convergence}

\begin{definition}[Strong convergence]
\label{def:strong}
A discretisation scheme is said to converge \emph{strongly} to the
solution of~\eqref{eq:sde} with order $\gamma > 0$ if there exist
constants $C > 0$ and $h_0 > 0$ such that, when the scheme is run on
the Brownian path of $X$,
\begin{equation}
\label{eq:strong-defn}
\E\!\left[\, |X_T - \hat X_N| \,\right] \;\leq\; C\, h^\gamma
\qquad \text{for all } h \in (0, h_0].
\end{equation}
\end{definition}

The expectation is taken over the (shared) Brownian path. The
definition uses $L^1$ norm, which is the most common choice; some
references use $L^2$ instead, with a corresponding adjustment of the
constant $C$ but the same order $\gamma$ for all the schemes we will
consider.

The strong error~\eqref{eq:strong-defn} measures average pathwise
deviation. It is the relevant quantity whenever a single trajectory of
the discretised process is meant to stand in for a single trajectory
of the true process: in pathwise sensitivity analysis (greeks via the
pathwise method, see Block~4), in optimal control problems where the
control is adapted to the simulated path, and in any context where
two related quantities must be evaluated on the \emph{same} simulated
path.

\subsection{Weak convergence}

\begin{definition}[Weak convergence]
\label{def:weak}
A discretisation scheme is said to converge \emph{weakly} to the
solution of~\eqref{eq:sde} with order $\beta > 0$ if for every test
function $f$ in a specified class of smooth functions (typically
polynomials of degree $\leq 2 \beta$ with bounded derivatives, or
$C_b^{2\beta}$), there exist constants $C_f > 0$ and $h_0 > 0$ such
that
\begin{equation}
\label{eq:weak-defn}
\big| \E[f(X_T)] - \E[f(\hat X_N)] \big|
\;\leq\; C_f\, h^\beta
\qquad \text{for all } h \in (0, h_0].
\end{equation}
\end{definition}

The weak error compares expectations of test functions, which are
averages over all paths. No coupling between the discretised and the
exact processes is required by the definition: $\E[f(\hat X_N)]$
depends only on the distribution of $\hat X_N$, regardless of how
$\hat X_N$ is constructed.

The weak error~\eqref{eq:weak-defn} is the relevant quantity whenever
the Monte Carlo estimator is used to compute an expectation: option
pricing under risk-neutral measure, expected utility, the partition
function in physics, and so on. In each of these settings, what is
estimated by the simulation is an integral against the law of $X_T$,
and only the marginal distribution matters.

\subsection{Why the two are different}

The two notions are formally different because they probe different
features of the joint law of $(X_T, \hat X_N)$. Strong convergence
sees the dependence structure between $X$ and $\hat X$; weak
convergence sees only their marginal distributions. The following
elementary example illustrates how dramatic the difference can be.

\begin{example}[Independent Gaussians]
\label{ex:indep-gauss}
Let $X \sim \mathcal{N}(0, 1)$ and $\hat X \sim \mathcal{N}(0, 1)$ be
independent. Then $X$ and $\hat X$ have \emph{identical} marginal
distributions, so $\E[f(X)] = \E[f(\hat X)]$ for every $f$ and the
``weak error'' is zero. But the pathwise distance is large:
$\E[|X - \hat X|] = \E[|X|] + \E[|\hat X|]$ by independence of
$X$ and $\hat X$ ... no, more carefully,
$\E[|X - \hat X|] = \sqrt{2/\pi} \cdot \sqrt{2} = 2/\sqrt{\pi} \approx 1.13$,
which is firmly bounded away from zero.
\end{example}

The example is a degenerate case (no time step, no SDE), but it
captures the principle: the strong and weak notions can disagree
arbitrarily on objects with identical marginals. For SDE schemes the
difference is more subtle but operates in the same direction: weak
convergence exploits cancellation across paths that strong
convergence cannot see.

A useful one-way relation between the two does hold, however. Suppose
$X_T$ and $\hat X_N$ are coupled (run on the same Brownian path), and
$f$ is Lipschitz with constant $L_f$. Then by Jensen's inequality,
\begin{equation}
\label{eq:weak-lip-strong}
\big| \E[f(X_T)] - \E[f(\hat X_N)] \big|
\;=\; \big| \E[f(X_T) - f(\hat X_N)] \big|
\;\leq\; \E[|f(X_T) - f(\hat X_N)|]
\;\leq\; L_f\, \E[|X_T - \hat X_N|].
\end{equation}
So strong order $\gamma$ implies weak order $\geq \gamma$ for
Lipschitz test functions. The converse fails: weak order can strictly
exceed strong order, and we will see this happen for Euler--Maruyama
(strong $\tfrac{1}{2}$, weak $1$).

\subsection{Which one matters in this project}

For European option pricing under a risk-neutral measure, the quantity
of interest is $\E[e^{-rT} \Phi(X_T)]$ for a payoff $\Phi$. This is
weak: marginal distributions only. \emph{Strong convergence is
irrelevant for the European pricing of Block~1.2.1 and~1.2.2.} What
we require, and what we will measure empirically, is weak order.

For path-dependent pricing in Block~5, the relevant question becomes
the weak error on functionals of the entire path
$\E[\Phi(X_{[0,T]})]$, not just on $\E[\Phi(X_T)]$. The marginals at
the grid points are not enough; in particular, for barrier options the
discretised supremum has its own bias of order $\sqrt{h}$ that adds to
any scheme's marginal error, and a Brownian-bridge correction will be
deployed. Strong convergence remains irrelevant to barrier pricing
\emph{per se}.

For pathwise sensitivity analysis in Block~4, by contrast, strong
convergence becomes essential: the pathwise estimator differentiates
the simulated path with respect to a parameter, and that
differentiation requires the path to be a faithful pathwise
approximation of the true process.

We therefore care about both notions across the project, but in
different blocks. For Block~1.2 specifically the operational
conclusion is that we will measure both orders empirically (the strong
order is interesting because it is the surprising one: $\tfrac{1}{2}$
for Euler), but only the weak order is the metric that justifies
using a given scheme for European pricing.

\section{The It\^o-Taylor expansion}
\label{sec:ito-taylor}

This section develops the device by which discretisation schemes for
SDEs are derived systematically. The mechanics generalise the
deterministic Taylor expansion, with the role of the chain rule
played by It\^o's formula. The treatment is detailed because the
expansion is the principal source of conceptual novelty in this block.

\subsection{Reminder: Taylor expansion for ODEs}

Consider the deterministic ODE $\dot x = a(x)$. The Taylor expansion
of $x(t + h)$ around $x(t)$ to second order reads
\begin{equation}
x(t + h) \;=\; x(t) + h\, \dot x(t) + \tfrac{1}{2} h^2 \ddot x(t) + O(h^3).
\end{equation}
Using the ODE itself, $\dot x = a(x)$ and $\ddot x = a'(x)\, \dot x =
a'(x)\, a(x)$. Substituting,
\begin{equation}
\label{eq:ode-taylor}
x(t + h) \;=\; x(t) + h\, a(x(t)) + \tfrac{1}{2} h^2 a'(x(t))\, a(x(t))
+ O(h^3).
\end{equation}
The explicit Euler scheme for ODEs is precisely the truncation of
\eqref{eq:ode-taylor} after the first-order term:
$x_{n+1} = x_n + h\, a(x_n)$. Higher-order schemes (improved Euler,
Runge--Kutta) keep more terms.

This logic extends to SDEs but with two complications: (i) the
solution $X_t$ of~\eqref{eq:sde} is not differentiable in $t$, so we
cannot literally Taylor-expand in time; (ii) the chain rule for
$f(X_t)$ is It\^o's formula, which carries a quadratic-variation term
absent in the ODE case. The first difficulty is overcome by working
with the integral form of the SDE, the second by accepting It\^o's
formula as the SDE chain rule.

\subsection{It\^o's formula and the generator}
\label{sec:ito-and-generator}

Recall It\^o's formula for a $C^2$ function $f$:
\begin{equation}
\label{eq:itos-formula}
\dd f(X_t) \;=\;
\Lop f(X_t) \dd t \;+\; f'(X_t)\, b(X_t) \dd W_t,
\end{equation}
where the \emph{It\^o generator} $\Lop$ is the operator
\begin{equation}
\label{eq:generator}
\Lop f(x) \;\eqdef\; a(x) f'(x) + \tfrac{1}{2} b^2(x) f''(x).
\end{equation}
Integrating~\eqref{eq:itos-formula} from $t$ to $t + h$,
\begin{equation}
\label{eq:itos-integral}
f(X_{t+h}) \;=\; f(X_t) + \int_t^{t+h} \Lop f(X_s) \dd s
\;+\; \int_t^{t+h} f'(X_s)\, b(X_s) \dd W_s.
\end{equation}

Two interpretations of $\Lop$ are useful. As a differential operator,
$\Lop$ is the second-order operator that gives the drift of $f(X)$
under~\eqref{eq:sde}; the second-derivative term comes from the
quadratic variation of $W$. As an analogue of $a \cdot f'$ in the
deterministic case, $\Lop$ plays the role of the ``time derivative''
operator for SDEs: it is the Markov generator of the diffusion $X$,
and $\Lop f(x) = \lim_{h \to 0} h^{-1}\, (\E[f(X_{t+h}) | X_t = x] -
f(x))$.

The integral identity~\eqref{eq:itos-integral} is the SDE analogue of
the fundamental theorem of calculus. The It\^o-Taylor expansion is what
results from \emph{iterating} this identity.

\subsection{First iteration: the SDE in integral form}

Apply~\eqref{eq:itos-integral} with $f(x) = x$. Since $f' = 1$ and
$f'' = 0$, the generator gives $\Lop f = a$, and
$f' \cdot b = b$. The result is
\begin{equation}
\label{eq:sde-integral}
X_{t+h} \;=\; X_t + \int_t^{t+h} a(X_s) \dd s
\;+\; \int_t^{t+h} b(X_s) \dd W_s.
\end{equation}
This is just the SDE~\eqref{eq:sde} in integral form, no surprise.
What it gives us, however, is a starting point: $X_{t+h}$ is expressed
in terms of $X_t$ plus two integrals over $[t, t+h]$ involving
$a(X_s)$ and $b(X_s)$. If we could replace $a(X_s)$ and $b(X_s)$ by
$a(X_t)$ and $b(X_t)$ at the cost of a small error, we would have
\begin{equation}
X_{t+h} \;\approx\; X_t + a(X_t) h + b(X_t) \Delta W,
\end{equation}
which is the Euler--Maruyama scheme. The question is what ``small
error'' means and how to quantify it. We answer this by applying
\eqref{eq:itos-integral} again, this time to $a$ and to $b$.

\subsection{Second iteration: expanding the integrands}
\label{sec:second-iteration}

Apply~\eqref{eq:itos-integral} with $f = a$:
\begin{equation}
\label{eq:expand-a}
a(X_s) \;=\; a(X_t) + \int_t^s \Lop a(X_u) \dd u
\;+\; \int_t^s a'(X_u)\, b(X_u) \dd W_u
\qquad \text{for } s \in [t, t+h],
\end{equation}
and similarly with $f = b$:
\begin{equation}
\label{eq:expand-b}
b(X_s) \;=\; b(X_t) + \int_t^s \Lop b(X_u) \dd u
\;+\; \int_t^s b'(X_u)\, b(X_u) \dd W_u.
\end{equation}
Substitute~\eqref{eq:expand-a} and~\eqref{eq:expand-b}
into~\eqref{eq:sde-integral}. The first integral becomes
\begin{equation}
\int_t^{t+h} a(X_s) \dd s
\;=\; a(X_t) \int_t^{t+h} \dd s
\;+\; \int_t^{t+h}\!\!\int_t^s \Lop a(X_u) \dd u \dd s
\;+\; \int_t^{t+h}\!\!\int_t^s a'(X_u) b(X_u) \dd W_u \dd s.
\end{equation}
The second integral becomes
\begin{equation}
\int_t^{t+h} b(X_s) \dd W_s
\;=\; b(X_t) \int_t^{t+h} \dd W_s
\;+\; \int_t^{t+h}\!\!\int_t^s \Lop b(X_u) \dd u \dd W_s
\;+\; \int_t^{t+h}\!\!\int_t^s b'(X_u) b(X_u) \dd W_u \dd W_s.
\end{equation}
Combining,
\begin{equation}
\label{eq:two-iterations}
X_{t+h} \;=\; X_t \;+\; a(X_t) h \;+\; b(X_t) \Delta W
\;+\; b(X_t) b'(X_t)\, I_{(1,1)}
\;+\; R,
\end{equation}
where:
\begin{itemize}
\item $\Delta W = W_{t+h} - W_t = \int_t^{t+h} \dd W_s$,
\item $I_{(1,1)} \eqdef \int_t^{t+h}\!\!\int_t^s \dd W_u \dd W_s$ is
the iterated It\^o integral over two Brownian increments,
\item $R$ collects all remaining terms: the drift--drift double
integral $\int\!\int \Lop a$, the drift--diffusion mixed
$\int\!\int a' b$, and so on, plus the leftover integrand factors
that we replaced by their values at $t$.
\end{itemize}
We have deliberately written~\eqref{eq:two-iterations} so that the
$b\, b' \cdot I_{(1,1)}$ term is isolated, because it is the term
that distinguishes Milstein from Euler.

\subsection{Closed form for the iterated integral}
\label{sec:iterated-integral}

The iterated It\^o integral $I_{(1,1)}$ admits a closed form. Apply
It\^o's formula to $f(W_t) = \tfrac{1}{2} W_t^2$:
\begin{equation}
\dd \!\left(\tfrac{1}{2} W_t^2\right) \;=\; W_t \dd W_t + \tfrac{1}{2} \dd t,
\end{equation}
which integrated from $t$ to $t+h$ gives
\begin{equation}
\tfrac{1}{2}(W_{t+h}^2 - W_t^2) \;=\; \int_t^{t+h} W_s \dd W_s
\;+\; \tfrac{1}{2} h.
\end{equation}
Subtract $W_t \int_t^{t+h} \dd W_s = W_t \Delta W$ from both sides
and use $W_s - W_t = \int_t^s \dd W_u$ for $s \in [t, t+h]$:
\begin{equation}
\int_t^{t+h} (W_s - W_t) \dd W_s
\;=\; \tfrac{1}{2}((W_{t+h} - W_t)^2 - h)
\;=\; \tfrac{1}{2}\big((\Delta W)^2 - h\big).
\end{equation}
The left-hand side is exactly $I_{(1,1)}$ (using
$W_s - W_t = \int_t^s \dd W_u$ and Fubini for stochastic integrals,
which holds in this It\^o setting). Therefore
\begin{equation}
\label{eq:I11}
I_{(1,1)} \;=\; \tfrac{1}{2}\big( (\Delta W)^2 - h \big).
\end{equation}

This identity is the key technical ingredient of the Milstein scheme:
it expresses the iterated integral, which is in principle a
$\sigma$-algebra-rich object built from the Brownian increment plus
its variance structure, in terms of the Brownian increment alone. We
do not need to simulate $I_{(1,1)}$ separately; given $\Delta W \sim
\mathcal{N}(0, h)$, we already know $I_{(1,1)}$ exactly.

\begin{remark}[Multidimensional case]
For SDEs driven by several independent Brownian motions $W^{(1)},
\dots, W^{(d)}$, the analogue of~\eqref{eq:I11} for cross-iterated
integrals such as $\int\!\int \dd W^{(i)} \dd W^{(j)}$ with $i \neq j$
\emph{does not} have a closed form in terms of the increments.
Simulating these cross-iterated integrals exactly requires extra work
(or a Brownian bridge construction), and this is the reason the
Milstein scheme in dimension $d \geq 2$ is significantly more
demanding than in dimension $1$. Phase~2 stays in the scalar setting,
so~\eqref{eq:I11} is sufficient.
\end{remark}

\subsection{The expansion to second order}

Substituting~\eqref{eq:I11} into~\eqref{eq:two-iterations}:
\begin{equation}
\label{eq:taylor-order-two}
X_{t+h} \;=\; X_t + a(X_t) h + b(X_t) \Delta W
\;+\; \tfrac{1}{2} b(X_t) b'(X_t)\, \big[(\Delta W)^2 - h\big]
\;+\; R.
\end{equation}
The remainder $R$ contains all the terms that have not been written
explicitly: the double drift integral, the mixed integrals, and the
contributions from the integrands $\Lop a, a' b, \Lop b, b' b$
evaluated at points $X_u$ rather than $X_t$. Each of these terms is,
on closer inspection, of higher order in $h$ than the explicit terms:
the double drift integral is $O(h^2)$, the mixed drift--diffusion
integrals are $O(h^{3/2})$ in $L^1$, and so on. The next subsection
makes this precise informally.

\subsection{Sizes of stochastic increments}
\label{sec:sizes}

A useful heuristic for predicting which terms matter at which order is
the rule of thumb
\begin{equation}
\Delta W \;\sim\; \sqrt{h},
\qquad h \;\sim\; h.
\end{equation}
That is, the random increment $\Delta W$ has standard deviation
$\sqrt{h}$, so its ``typical size'' in any norm sensitive to it is
$\sqrt{h}$, while a deterministic increment of size $h$ stays $h$.
Combinations:
\begin{itemize}
\item $(\Delta W)^2 \sim h$ (since $\Var(\Delta W) = h$),
\item $\Delta W \cdot h \sim h^{3/2}$,
\item $(\Delta W)^3 \sim h^{3/2}$ (third moment of $\mathcal{N}(0, h)$
is $0$, but typical magnitudes are $h^{3/2}$),
\item double iterated integrals like
$\int\!\int \dd W \dd W \sim h$ (this is precisely $I_{(1,1)}$ from
\eqref{eq:I11}),
\item double iterated integrals like $\int\!\int \dd s \dd s \sim h^2$.
\end{itemize}

This calculus is informal but accurate enough to predict the orders
of all schemes we will encounter. With it, the structure of
\eqref{eq:taylor-order-two} becomes:
\begin{itemize}
\item $a(X_t)\, h$: order $h$.
\item $b(X_t)\, \Delta W$: order $h^{1/2}$.
\item $\tfrac{1}{2} b\, b'\, [(\Delta W)^2 - h]$: order $h$.
\item $R$ (terms not written): order $h^{3/2}$ and higher.
\end{itemize}

\section{The Euler--Maruyama scheme}
\label{sec:euler}

\subsection{Derivation as first-order truncation}

The Euler--Maruyama scheme is obtained by truncating
\eqref{eq:taylor-order-two} after the $\Delta W$ term: keep $h$ and
$h^{1/2}$, drop $h$ from $I_{(1,1)}$ and everything beyond. The
recursion is
\begin{equation}
\label{eq:euler}
\hat X_{n+1} \;=\; \hat X_n
\;+\; a(\hat X_n)\, h
\;+\; b(\hat X_n)\, \Delta W_n,
\qquad \hat X_0 = x_0.
\end{equation}
This is the natural extension to SDEs of the explicit Euler scheme
for ODEs: replace $a$ by $b\, dW$ where appropriate and add the
stochastic increment.

The truncation drops, at each step, a quantity that
\eqref{eq:taylor-order-two} reveals to be of size $h$ (the
$I_{(1,1)}$-related term) plus higher-order remainders.

\subsection{Convergence orders}

\begin{theorem}[Convergence of Euler--Maruyama]
\label{thm:euler-conv}
Suppose $a, b \in C^2(\R)$ with $a, b, a', b'$ Lipschitz and bounded.
The Euler--Maruyama scheme~\eqref{eq:euler} converges to the solution
of~\eqref{eq:sde} with strong order $\gamma = \tfrac{1}{2}$ and weak
order $\beta = 1$.
\end{theorem}

The proof is technical and is given in Kloeden--Platen,
\S\,10.2 (strong) and \S\,14.5 (weak). We do not reproduce it. What
we do is explain, at the level of intuition, why the orders are
$\tfrac{1}{2}$ and $1$ respectively.

\paragraph{Why strong order $\tfrac{1}{2}$.} The local truncation
error of one Euler step is the dropped term
$\tfrac{1}{2} b b' [(\Delta W)^2 - h] + O(h^{3/2})$.
In $L^1$, $|(\Delta W)^2 - h|$ has expectation
$\E[|(\Delta W)^2 - h|] \sim h$ (this is a moment of a centred chi-square), so the local
strong error is $O(h)$. There are $N = T/h$ steps and the errors
combine in $L^1$ via a Gronwall-type argument that is essentially a
square-root law (the variance accumulates linearly, not the standard
deviation). The global strong error is thus $\sqrt{N \cdot h^2} =
\sqrt{T h}$, that is, $O(h^{1/2})$.

\paragraph{Why weak order $1$.} The local truncation error has
expectation $\tfrac{1}{2} \E[b b' \cdot ((\Delta W)^2 - h)]$, which
\emph{vanishes} because $\E[(\Delta W)^2 - h] = 0$. The leading
contribution to the weak error therefore comes from the next term in
the expansion, which is $O(h^{3/2})$ but again has zero mean to
leading order; what survives are second-moment terms of size $O(h^2)$
locally, accumulating to $O(h)$ globally. Hence weak order $1$.

The asymmetry $\gamma = \tfrac{1}{2}, \beta = 1$ is therefore not an
accident: it reflects the fact that the dominant ``missing'' term
$\tfrac{1}{2} b b' [(\Delta W)^2 - h]$ is path-dependent (so seen by
$L^1$) but mean-zero (so invisible to $\E[\cdot]$).

\subsection{What the order means in practice}

Suppose we want to halve the strong error: we must \emph{quarter} the
step size $h$ (since strong error $\sim h^{1/2}$). The number of steps
$N = T/h$ quadruples and so does the simulation cost (per path). To
halve the weak error: we must halve the step size; cost only doubles.

For European pricing, where weak order is what matters, Euler is
respectable. For pathwise applications, the strong order
$\tfrac{1}{2}$ becomes a real bottleneck and is the motivation for
moving to Milstein.

\subsection{Application to GBM}

For geometric Brownian motion, $a(x) = r x$ and $b(x) = \sigma x$, so
\eqref{eq:euler} becomes
\begin{equation}
\label{eq:euler-gbm}
\hat S_{n+1} \;=\; \hat S_n + r\, \hat S_n\, h
\;+\; \sigma\, \hat S_n\, \Delta W_n
\;=\; \hat S_n \big(1 + r h + \sigma \Delta W_n\big).
\end{equation}
The recursion is multiplicative, which is convenient. A nuisance,
specific to GBM under Euler: there is no a priori guarantee that
$\hat S_n$ remains positive: if a $\Delta W_n$ is negative enough that
$1 + r h + \sigma \Delta W_n < 0$, the simulated price becomes
negative or zero, which is unphysical. The probability of this for
realistic parameters and reasonable $h$ is small but nonzero, and the
log-Euler scheme (apply Euler in $\log S$, then exponentiate)
sidesteps the issue. We will mention this in the implementation but
not adopt it: the standard Euler is the canonical scheme and exhibits
the canonical convergence behaviour we want to validate.

\section{The Milstein scheme}
\label{sec:milstein}

\subsection{Derivation as second-order truncation}

The Milstein scheme is obtained by keeping the $I_{(1,1)}$ term in
\eqref{eq:taylor-order-two}, that is, truncating after order $h$:
\begin{equation}
\label{eq:milstein}
\hat X_{n+1} \;=\; \hat X_n
\;+\; a(\hat X_n)\, h
\;+\; b(\hat X_n)\, \Delta W_n
\;+\; \tfrac{1}{2} b(\hat X_n)\, b'(\hat X_n)\,
\big[(\Delta W_n)^2 - h\big],
\qquad \hat X_0 = x_0.
\end{equation}
The new term, called the Milstein correction, is the iterated It\^o
integral $I_{(1,1)}$ multiplied by $b\, b'$ evaluated at the current
state. By~\eqref{eq:I11} this term is computable directly from
$\Delta W_n$ and $h$, with no extra random variables. The cost of
Milstein over Euler is the evaluation of $b'(\hat X_n)$ and a few
extra arithmetic operations.

\subsection{Convergence orders}

\begin{theorem}[Convergence of Milstein]
\label{thm:milstein-conv}
Suppose $a, b \in C^2(\R)$ with $a, b, a', b', a'', b''$ Lipschitz
and bounded. The Milstein scheme~\eqref{eq:milstein} converges to the
solution of~\eqref{eq:sde} with strong order $\gamma = 1$ and weak
order $\beta = 1$.
\end{theorem}

\paragraph{Why strong order $1$.} By including the
$I_{(1,1)}$ term, Milstein eliminates the $O(h)$ local strong error
of Euler. The next remainder is $O(h^{3/2})$, which accumulates over
$N = T/h$ steps to $O(h)$ globally. Hence strong order $1$.

\paragraph{Why weak order remains $1$.} Adding a mean-zero correction
to Euler does not improve the weak error: the leading non-zero
expectation term is unchanged. To get weak order $> 1$, one needs to
correct further moments of the increment; this requires schemes like
the order-$2$ weak scheme of Platen, which is genuinely more involved
and beyond the scope of this block.

\subsection{Application to GBM}

For GBM, $b(x) = \sigma x$, so $b'(x) = \sigma$ and
$b(x) b'(x) = \sigma^2 x$. The recursion becomes
\begin{equation}
\label{eq:milstein-gbm}
\hat S_{n+1} \;=\; \hat S_n + r \hat S_n h + \sigma \hat S_n \Delta W_n
+ \tfrac{1}{2} \sigma^2 \hat S_n \big[(\Delta W_n)^2 - h\big].
\end{equation}
Again multiplicative, again with a positivity issue (although the
correction term raises the lower bound: the squared term
$(\Delta W_n)^2$ is always nonnegative, so the correction is bounded
below by $-\tfrac{1}{2} \sigma^2 \hat S_n h$, and the worst-case loss
per step is mitigated). We will report empirically what fraction of
paths, if any, develop nonpositive prices in our test parameters.

\section{Practical considerations and which scheme to use}
\label{sec:practical}

The textbook recommendation can be misleading and deserves an
explicit unpacking. Phrases like ``Milstein is more accurate than
Euler'' are true in the strong sense and \emph{false} in the weak
sense, where the two schemes are tied at order $1$.

\subsection{For European option pricing}

We compute $\E[e^{-rT} \Phi(S_T)]$. This is a marginal expectation,
the weak setting. Both Euler and Milstein have weak order $1$.
Milstein costs more per step (one extra multiplication and one
derivative evaluation; the latter can be analytically symbolic and
hence ``free'' if the model coefficients are simple). The only
benefit of Milstein over Euler in this setting is a smaller constant
in front of the $h$ asymptote, not an improved order.

In practice for pricing under GBM (and for many local volatility
models), the constant improvement is too small to justify the extra
machinery: Euler with a slightly smaller $h$ achieves the same weak
error at lower or comparable cost. Block~1.2.1 will demonstrate
empirically that for European pricing under GBM, Euler at
$h = T/200$ and Milstein at $h = T/100$ (matched cost) deliver
comparable accuracy.

\subsection{For pathwise sensitivity (Block~4)}

The pathwise method estimates greeks by differentiating the simulated
path with respect to a parameter, then averaging. For the
differentiation to be meaningful, the simulated path must be a
faithful pathwise approximation, which is the strong setting.
Milstein's strong order $1$ versus Euler's $\tfrac{1}{2}$ becomes a
genuine advantage: Milstein needs $\sqrt{h}$ as many steps as Euler
to attain the same strong error, an asymptotic factor that grows
without bound as $h \to 0$.

\subsection{For path-dependent payoffs (Block~5)}

The expectation is now over the path, $\E[\Phi(S_{[0,T]})]$. Two
sources of error: the marginal scheme error (weak, both schemes
$O(h)$) and the discretisation of the payoff functional itself
(barriers, suprema). The latter dominates for barrier options at
$O(\sqrt{h})$, regardless of scheme. So Euler is sufficient here,
and the relevant improvement is in the payoff treatment (Brownian
bridge correction), not in the scheme.

\subsection{Summary}

The verdict is unambiguous: \emph{for the European pricing of
Block~1.2, Euler--Maruyama is the right choice and Milstein is
implemented for completeness and as a comparison baseline rather than
because it offers any operational advantage}. Block~1.2.1 will
implement Euler as the working pricer; Block~1.2.2 will implement
Milstein as a side-by-side benchmark. The empirical convergence
study at the heart of validation will measure both strong and weak
orders for both schemes, confirming the predictions of
Theorems~\ref{thm:euler-conv} and~\ref{thm:milstein-conv}.

\section{Roadmap for Blocks 1.2.1 and 1.2.2}
\label{sec:roadmap}

Block~1.2.1 will implement Euler--Maruyama for GBM and validate it
empirically. The validation has three components:
\begin{enumerate}
\item \textbf{Strong-order verification}: simulate $\hat S_T$ on the
\emph{same} Brownian path as the exact $S_T$ (the latter computed by
the closed-form formula evaluated at $W_T = \sum_n \Delta W_n$),
measure $\E[|S_T - \hat S_T|]$ across many paths, plot in log-log
against $h$, fit a line, expect slope $\tfrac{1}{2}$.
\item \textbf{Weak-order verification}: with $\Phi(s) = (s - K)^+
e^{-rT}$, measure $|\E[\Phi(S_T)] - \E[\Phi(\hat S_T)]|$ across many
paths and many $h$, plot in log-log, expect slope $1$.
\item \textbf{Sanity checks}: containment of the BS price by the MC
CI for moderate $h$, monotonicities, reduction to Block~1.1's exact
sampler in the limit $h \to 0$ at fixed $n$.
\end{enumerate}
The first two are direct empirical confirmations of the convergence
theorems; the third anchors the implementation in the existing
infrastructure.

Block~1.2.2 will implement Milstein under the same protocol. The
strong order should change from $\tfrac{1}{2}$ to $1$; the weak order
should remain $1$, matching Euler. This is the canonical demonstration
of the asymmetry between the two notions of convergence.

Architecturally, both schemes will be added to the existing
\texttt{quantlib.monte\_carlo} (Python) and \texttt{quant/monte\_carlo}
(C++) modules, alongside the exact GBM sampler from Block~1.1. The
\texttt{mc\_estimator} reducer is unchanged. The pricer wrappers
(\texttt{mc\_european\_call\_euler}, \texttt{mc\_european\_call\_milstein})
will be thin orchestrators on top of new path samplers.

\section{Bibliographic notes}
\label{sec:biblio}

The canonical reference for the numerical analysis of SDEs is
Kloeden \& Platen~\cite{KloedenPlaten}, especially Chapters~5
(stochastic Taylor expansions), 9--10 (Euler--Maruyama and strong
convergence), and 14--15 (weak convergence and Milstein-type
schemes). The treatment is comprehensive and self-contained but
demanding; the relevant theorems are stated with multiple regularity
hypotheses that we have stylised in Theorems~\ref{thm:euler-conv}
and~\ref{thm:milstein-conv}. For a reader-friendlier introduction,
Higham~\cite{Higham2001} is a classic survey article that develops
exactly the empirical convergence demonstration we will reproduce in
the implementation. Glasserman~\cite[Chapter~6]{Glasserman} treats
the topic from the financial Monte Carlo angle and includes
discussion of the practical relevance of weak versus strong
convergence in the contexts we will encounter in subsequent blocks.

\begin{thebibliography}{99}

\bibitem{Glasserman}
P.\ Glasserman,
\emph{Monte Carlo Methods in Financial Engineering},
Springer, 2004.

\bibitem{Higham2001}
D.\ J.\ Higham,
\emph{An algorithmic introduction to numerical simulation of
stochastic differential equations},
SIAM Review, 43(3):525--546, 2001.

\bibitem{KloedenPlaten}
P.\ E.\ Kloeden and E.\ Platen,
\emph{Numerical Solution of Stochastic Differential Equations},
Springer, 1992.

\bibitem{Etheridge}
A.\ Etheridge,
\emph{A Course in Financial Calculus},
Cambridge University Press, 2002.

\end{thebibliography}

\end{document}
